\subsection{Exponential Memory Decay and the Origin of the IR-Safe Sector}

The repeated emergence of \(n\simeq 1\) in the observationally viable SBC/F3 region suggests that the infrared correction may originate from a first-order memory-decay process.

Let \(M(\tau)\) denote an effective cosmological memory amplitude, where \(\tau\) parametrizes the causal depth or accumulated observational scale. The simplest stable decay law is

\begin{equation}
\frac{dM}{d\tau}=-\lambda M,
\end{equation}

with solution

\begin{equation}
M(\tau)=M_0 e^{-\lambda \tau}.
\end{equation}

If the relevant infrared causal depth scales approximately as

\begin{equation}
\tau \propto (1+z),
\end{equation}

then the effective SBC/F3 response becomes

\begin{equation}
F_{\rm SBC}(z)
\simeq
-\epsilon \exp[-k_0(1+z)].
\end{equation}

This is the leading-order case of the phenomenological form

\begin{equation}
F_{\rm SBC}(z)
=
-\epsilon
\exp\left[-\left(k_0(1+z)\right)^n\right],
\end{equation}

with

\begin{equation}
n\simeq 1.
\end{equation}

Thus, the observational preference for \(n\simeq 1\) can be interpreted as the signature of a first-order exponentially screened memory response.

In this interpretation, \(\epsilon\) measures the amplitude of the effective memory imprint, while \(k_0\) controls the infrared screening scale. The preferred region,

\begin{equation}
\epsilon\simeq 0.06,
\qquad
k_0\simeq 0.015-0.022,
\qquad
n\simeq 1,
\end{equation}

corresponds to a weak, slowly decaying, large-scale memory sector.

This provides a physical rationale for why the SBC/F3 deformation is observationally viable: it is sufficiently weak to preserve the \(\Lambda\)CDM background and high-\(\ell\) acoustic structure, while remaining large enough to generate controlled infrared perturbative effects.
